\documentclass[11pt]{article}
\usepackage[T1]{fontenc}
\usepackage[utf8]{inputenc}
\usepackage{lmodern}
\usepackage[margin=1in]{geometry}
\usepackage{microtype}
\usepackage{amsmath,amssymb,amsthm,mathtools}
\usepackage{booktabs,longtable,array}
\usepackage{enumitem}
\usepackage{xcolor}
\usepackage{hyperref}
\usepackage[nameinlink,capitalise]{cleveref}
\usepackage{siunitx}
\usepackage{listings}
\usepackage{fancyhdr}
\usepackage{setspace}

\definecolor{deepblue}{RGB}{24,55,95}
\definecolor{deepgreen}{RGB}{28,100,70}
\definecolor{deepred}{RGB}{145,45,45}
\definecolor{softgray}{RGB}{245,247,250}
\hypersetup{colorlinks=true,linkcolor=deepblue,citecolor=deepgreen,urlcolor=deepblue,pdftitle={Positive Interpolation Rigidity and New Routes toward the Alaoglu--Erdos Conjecture}}
\pagestyle{fancy}
\fancyhf{}
\lhead{Alaoglu--Erd\H{o}s progress report}
\rhead{22 August 2026}
\cfoot{\thepage}
\setlength{\headheight}{14pt}
\setstretch{1.04}
\setlist{itemsep=0.25em,topsep=0.4em}
\sisetup{scientific-notation=true,round-mode=places,round-precision=7}

\newtheorem{theorem}{Theorem}[section]
\newtheorem{proposition}[theorem]{Proposition}
\newtheorem{lemma}[theorem]{Lemma}
\newtheorem{corollary}[theorem]{Corollary}
\newtheorem{conjecture}[theorem]{Conjecture}
\newtheorem{problem}[theorem]{Problem}
\theoremstyle{definition}
\newtheorem{definition}[theorem]{Definition}
\newtheorem{remark}[theorem]{Remark}
\newtheorem{experiment}[theorem]{Experiment}

\newcommand{\N}{\mathbb N}
\newcommand{\Z}{\mathbb Z}
\newcommand{\Q}{\mathbb Q}
\newcommand{\R}{\mathbb R}
\newcommand{\C}{\mathbb C}
\newcommand{\E}{\mathbb E}
\newcommand{\Var}{\operatorname{Var}}
\newcommand{\Cov}{\operatorname{Cov}}
\newcommand{\supp}{\operatorname{supp}}
\newcommand{\dist}{\operatorname{dist}}
\newcommand{\vct}[1]{\boldsymbol{#1}}
\newcommand{\status}[2]{\par\medskip\noindent\colorbox{#1!10}{\parbox{0.965\linewidth}{\textbf{#2}}}\par\medskip}
\newcommand{\proved}{\status{deepgreen}{PROVED IN THIS REPORT. }}
\newcommand{\openstatus}{\status{deepred}{OPEN BOTTLENECK. }}
\newcommand{\conditional}{\status{deepblue}{CONDITIONAL RESEARCH TARGET. }}

\lstset{basicstyle=\ttfamily\small,breaklines=true,frame=single,backgroundcolor=\color{softgray},columns=fullflexible,keepspaces=true,showstringspaces=false}

\title{\textbf{Positive Interpolation Rigidity, Geometric Newton Expansions, and Integer-Cell Dynamics}\\[0.4em]
\large New progress toward the Alaoglu--Erd\H{o}s conjecture}
\author{Research companion report}
\date{22 August 2026}

\begin{document}
\maketitle

\begin{abstract}
Let $x\in\R$ and suppose that $a=2^x$ and $b=3^x$ are integers.  The Alaoglu--Erd\H{o}s conjecture asserts that $x$ is an integer.  This document is a companion to the supplied unified report, not a synopsis of it.  It deliberately avoids reproducing that report's determinant, transcendence, $p$-adic, moment-matrix, and formalization programs.  The principal new result is an exact four-node covariance identity: a power series with nonnegative coefficients that assumes the values $(1,a,b,ab)$ at $(1,2,3,6)$ must be a single monomial.  This gives a sharp rigidity theorem, a quantitative stability estimate, an explicit Farkas-type separating functional, and a weighted lower bound on the negative mass of every signed interpolant.

A second new line develops the Newton interpolation of $z^{\log_2 3}$ on the geometric grid $1,2,4,\dots$.  Its coefficients admit a closed product formula, have an eventually alternating sign, and produce a one-signed tail at every integer between consecutive powers of $2$.  Rewriting the terms as a basic hypergeometric series produces a canonical entire interpolant of $2^m\mapsto 3^m$ and isolates a concrete arithmetic comparison problem.  A third line packages every hypothetical counterexample into a positive integer cocycle over the irrational rotation $n\mapsto\{nx\}$; primitive counterexamples impose exact $2$-adic or $3$-adic signatures on that cocycle.

No complete proof is claimed.  The report identifies exactly what has been proved, why several tempting completions fail, and four narrowly formulated next targets.  Competitive announcements mentioned in the prompt are not used as mathematical premises; every result below is proved from stated assumptions.
\end{abstract}

\tableofcontents
\newpage

\section{Mandate, nonduplication, and status}

The supplied source is a large cumulative research report.  Its main families of ideas include interpolation determinants, transcendence conjectures, $p$-adic logarithms, auxiliary functions, matrix positivity, moment methods, approximation programs, and formalization plans.  Repeating those programs under new names would add little.  The present report therefore concentrates on three constructions that are structurally different:
\begin{enumerate}[label=(\roman*)]
\item a \emph{coefficientwise positive interpolation cone} and its exact covariance defect;
\item a \emph{geometric Newton basis} whose coefficients can be evaluated in closed form;
\item an \emph{integer-valued chord cocycle} attached to the fractional parts of $nx$.
\end{enumerate}
The elementary reductions in \cref{sec:reductions} are included only to make the document logically self-contained.

\proved The central positivity theorem, its measure-theoretic form, quantitative stability, the global separator, the signed-negativity budget, the geometric-Newton coefficient formula, and the integer cocycle are unconditional.

\openstatus The missing implication is not a technicality: integrality of $2^x$ and $3^x$ does not presently yield a positive integer-exponent representation.  The report turns that vague gap into explicit arithmetic subproblems, but does not close it.

\section{Baseline reductions and a primitive normalization}\label{sec:reductions}

Throughout, put
\[
 a=2^x,\qquad b=3^x,\qquad \alpha=\frac{\log 3}{\log 2}.
\]
Then $b=a^\alpha$ and $6^x=ab$.  Since a positive integer cannot lie strictly between $0$ and $1$, $x\ge 0$; the case $x=0$ is immediate.

\begin{proposition}[Rational and algebraic cases]\label{prop:basic-cases}
If $x\in\Q$ and $2^x\in\Z$, then $x\in\Z$.  If $x$ is algebraic and both $2^x$ and $3^x$ are algebraic, then $x\in\Z$.
\end{proposition}
\begin{proof}
Write $x=p/q$ in lowest terms with $q>0$.  If $2^{p/q}=a\in\Z$, then $a^q=2^p$; comparison of the $2$-adic valuations gives $q\mid p$, hence $q=1$.  For algebraic irrational $x$, the Gelfond--Schneider theorem says that every value of $2^x$ is transcendental, contradicting $2^x\in\Z$.  The remaining algebraic possibilities are rational and were just handled.
\end{proof}

Thus a counterexample would necessarily have $x$ transcendental.  It would also give the $2\times2$ algebraic array required to contradict the four exponentials conjecture: the two logarithms $\log2,\log3$ are linearly independent over $\Q$, as are $1,x$ when $x\notin\Q$, while all four exponentials $2,3,2^x,3^x$ would be algebraic.  The six exponentials theorem does not directly apply because no third independent logarithmic direction is available; this is precisely why the problem is delicate \cite{Waldschmidt2000,Lang1966}.

A small normalization is useful in every arithmetic approach.

\begin{proposition}[Primitive counterexample reduction]\label{prop:primitive}
Suppose $(x,a,b)$ is a nonintegral counterexample.  Let
\[
 r=v_2(a),\qquad s=v_3(b),\qquad t=\min(r,s).
\]
Then $y=x-t>0$ is again a nonintegral counterexample, with
\[
 2^y=\frac{a}{2^t}\in\Z,\qquad 3^y=\frac{b}{3^t}\in\Z,
\]
and at least one of $2^y$ and $3^y$ is a unit at its distinguished prime: either $v_2(2^y)=0$ or $v_3(3^y)=0$.
\end{proposition}
\begin{proof}
The displayed quotients are integers by the definition of $t$.  Since $2^t\le a=2^x$, one has $t\le x$, so $y\ge0$.  If $y=0$, then $x=t\in\Z$, contrary to hypothesis.  Subtracting an integer does not turn a noninteger into an integer.  Finally, one of $r-t$ and $s-t$ is zero.
\end{proof}

Accordingly, any minimal or normalized counterexample may be assumed to satisfy
\begin{equation}\label{eq:primitive-condition}
 a\ \text{odd}\qquad\text{or}\qquad 3\nmid b.
\end{equation}
This weak-looking condition later gives exact congruence information for the integer cocycle in \cref{sec:cocycle}.

\section{The four-node positivity identity}\label{sec:positivity}

The point of departure is to replace direct approximation of $z^x$ by a question about the cone generated by the integer monomials $1,z,z^2,\dots$.

\begin{theorem}[Four-node covariance identity]\label{thm:covariance-identity}
Let
\[
 F(z)=\sum_{n\ge0}c_n z^n,
\]
where $c_n\ge0$ and $F(6)<\infty$.  Then
\begin{align}
 F(1)F(6)-F(2)F(3)
 &=\sum_{0\le n<m}c_nc_m(2^m-2^n)(3^m-3^n) \label{eq:covariance-sum}\
 &\ge0.\nonumber
\end{align}
Equality holds if and only if at most one coefficient $c_n$ is nonzero.
\end{theorem}
\begin{proof}
All sums are absolutely convergent, so they may be rearranged.  Pair the ordered terms $(n,m)$ and $(m,n)$ in
\[
 \sum_{n,m}c_nc_m\bigl(6^m-2^n3^m\bigr).
\]
The diagonal terms vanish, while an unordered pair $n<m$ contributes
\[
 c_nc_m\bigl(6^m+6^n-2^n3^m-2^m3^n\bigr)
 =c_nc_m(2^m-2^n)(3^m-3^n).
\]
Every factor is nonnegative and is strictly positive for $n<m$.  Hence equality forces $c_nc_m=0$ for every distinct pair, which is equivalent to support of cardinality at most one.
\end{proof}

\begin{corollary}[Positive-interpolant rigidity]\label{cor:positive-rigidity}
Suppose $F(z)=\sum c_nz^n$ has nonnegative coefficients, $F(6)<\infty$, and
\[
 F(1)=1,\qquad F(2)=2^x,\qquad F(3)=3^x,\qquad F(6)=6^x.
\]
Then $x\in\N$ and $F(z)=z^x$.
\end{corollary}
\begin{proof}
The target values satisfy $F(1)F(6)=F(2)F(3)$.  By \cref{thm:covariance-identity}, $F(z)=cz^k$ for some $k\in\N$.  The value at $1$ gives $c=1$, and the value at $2$ gives $2^k=2^x$, hence $x=k$.
\end{proof}

The identity is not tied to the numbers $2$ and $3$.

\begin{theorem}[Measure form]\label{thm:measure-form}
Let $\mu$ be a probability measure on a totally ordered subset of $\R$, and let
\[
 M(u)=\int u^t\,d\mu(t)
\]
be finite at $u,v,uv$, where $u,v>1$.  Then
\begin{equation}\label{eq:measure-cov}
 M(uv)-M(u)M(v)
 =\frac12\iint (u^s-u^t)(v^s-v^t)\,d\mu(s)d\mu(t)\ge0.
\end{equation}
Equality holds if and only if $\mu$ is a point mass.
\end{theorem}
\begin{proof}
Expand the right-hand side and use symmetry.  Because $u^t$ and $v^t$ are strictly increasing functions of $t$, the integrand is nonnegative and vanishes exactly on the diagonal $s=t$.  Equality therefore implies that two independent samples from $\mu$ coincide almost surely, which is possible only for a point mass.
\end{proof}

In probabilistic language, if $N$ is integer-valued, then
\[
 M(6)-M(2)M(3)=\Cov(2^N,3^N).
\]
The two random variables are strictly comonotone, so zero covariance is a degeneracy certificate.  This is the shortest conceptual proof of \cref{cor:positive-rigidity}.

\subsection{Absolutely monotone interpolation}

A function analytic at the origin with nonnegative Taylor coefficients is absolutely monotone on every positive interval inside its disk of convergence.  Conversely, standard Bernstein-type representation theorems turn absolute monotonicity into coefficientwise positivity under mild analytic hypotheses.  Thus \cref{cor:positive-rigidity} may be read as follows.

\begin{corollary}\label{cor:absolute-monotone}
An analytic absolutely monotone interpolant through
\[
 (1,1),\quad (2,a),\quad (3,b),\quad (6,ab)
\]
can exist only when $(a,b)=(2^k,3^k)$ for an integer $k\ge0$.
\end{corollary}

This identifies a sharply delimited bridge: any construction that derives an absolutely monotone interpolant from the integrality of $a$ and $b$ would finish the conjecture immediately.  The later separator theorem explains why an exact construction cannot arise from a generic positivity-preserving interpolation routine; it has to use the arithmetic hypothesis in an essential way.

\section{Quantitative stability and denominator quantization}\label{sec:stability}

The exact identity has a useful stability form.  Let $(p_n)_{n\ge0}$ be a probability distribution and put $M(t)=\sum p_nt^n$ and
\[
 \Delta=M(6)-M(2)M(3).
\]

\begin{proposition}[Concentration from a small defect]\label{prop:stability}
One has
\begin{equation}\label{eq:stability}
 \Delta\ge 1-\sum_{n\ge0}p_n^2\ge 1-\sup_n p_n.
\end{equation}
Consequently, if $\Delta\le\varepsilon<1$, some atom has mass at least $1-\varepsilon$.
\end{proposition}
\begin{proof}
For $m>n$,
\[
 (2^m-2^n)(3^m-3^n)
 =6^n(2^{m-n}-1)(3^{m-n}-1)\ge2.
\]
Insert this in \eqref{eq:covariance-sum}:
\[
 \Delta\ge2\sum_{n<m}p_np_m=1-\sum_np_n^2.
\]
Since $p_n^2\le(\sup_jp_j)p_n$, the second inequality follows.
\end{proof}

There is also an exact arithmetic quantization.

\begin{corollary}[Rational-weight gap]\label{cor:rational-gap}
If every $p_n$ is an integral multiple of $1/D$, then $D^2\Delta$ is a nonnegative integer.  Hence either $\Delta=0$ or $\Delta\ge D^{-2}$.
\end{corollary}
\begin{proof}
Write $p_n=r_n/D$ with $r_n\in\N$.  Multiplying \eqref{eq:covariance-sum} by $D^2$ gives a sum of nonnegative integers.
\end{proof}

This gives a possible combinatorial completion criterion.  It is enough to construct a rational positive model whose covariance defect is strictly smaller than the square of its weight mesh.  The construction must be extraordinarily sharp: the target lies on the equality boundary, not in the interior of the moment cone.

\begin{remark}[Perturbed moments]
Suppose $M(2)=a+e_2$, $M(3)=b+e_3$, and $M(6)=ab+e_6$.  Then
\begin{equation}\label{eq:defect-errors}
 \Delta=e_6-be_2-ae_3-e_2e_3.
\end{equation}
Thus a rational positive approximation with controlled signed errors can be tested against \cref{cor:rational-gap} without evaluating a double sum.  This is useful for machine-assisted searches, but no construction currently reaches the required $D^{-2}$ scale.
\end{remark}

\section{The discrete moment cone and an explicit global separator}\label{sec:cone}

For $t\ge0$ define the four-coordinate moment curve
\[
 \vct v(t)=(1,2^t,3^t,6^t)\in\R^4,
\]
and let $\mathcal C$ be the convex cone generated by $\vct v(n)$ for $n\in\N$.  The target $\vct v(x)$ lies on the rank-one hypersurface $w=uv$, because $6^x=2^x3^x$.  By \cref{thm:covariance-identity}, every normalized point of $\mathcal C$ satisfies $w\ge uv$, with equality only at a single integer-exponent generator.

This nonlinear boundary statement has a remarkably simple linear dual certificate.

\begin{theorem}[Explicit Farkas separator]\label{thm:separator}
Let $k\in\N$ and define
\begin{align}
 h_k(t)
 &=(2^t-2^k)(3^t-3^{k+1})\label{eq:hk-factor}\\
 &=6^t-3^{k+1}2^t-2^k3^t+2^k3^{k+1}.\label{eq:hk-linear}
\end{align}
Then $h_k(n)\ge0$ for every integer $n\ge0$, while $h_k(t)<0$ for every $t\in(k,k+1)$.
\end{theorem}
\begin{proof}
For $n\le k$, both factors in \eqref{eq:hk-factor} are nonpositive; for $n\ge k+1$, both are nonnegative.  If $k<t<k+1$, the first factor is positive and the second negative.
\end{proof}

Thus the linear functional with coefficient vector
\[
 \lambda_k=(2^k3^{k+1},-3^{k+1},-2^k,1)
\]
is nonnegative on every integer-exponent generator and strictly negative at the noninteger point $\vct v(x)$.

\begin{corollary}[Integral separation margin]\label{cor:integral-margin}
If $a=2^x$ and $b=3^x$ are integers and $k<x<k+1$, then
\begin{equation}\label{eq:margin}
 h_k(x)=-(a-2^k)(3^{k+1}-b)\in\Z_{\le-1}.
\end{equation}
In particular, the separating margin is at least one before normalization.
\end{corollary}

\begin{corollary}[Sup-norm distance from the positive cone]\label{cor:distance}
Let $\vct y$ be any normalized positive mixture of the vectors $\vct v(n)$.  If each of its four coordinates differs from the corresponding coordinate of $\vct v(x)$ by at most $\varepsilon$, then
\[
 \varepsilon\ge \frac{1}{(2^k+1)(3^{k+1}+1)}.
\]
If the first coordinate is exactly $1$, the denominator may be replaced by $3^{k+1}+2^k+1$.
\end{corollary}
\begin{proof}
One has $\lambda_k\cdot\vct y\ge0$ and $\lambda_k\cdot\vct v(x)\le-1$.  Apply the triangle inequality to $\lambda_k\cdot(\vct y-\vct v(x))$.
\end{proof}

The same separator quantifies how much negativity a signed interpolant must contain.

\begin{theorem}[Weighted negativity budget]\label{thm:negativity-budget}
Suppose a finitely supported real sequence $(c_n)$ satisfies
\[
 \sum_nc_n\vct v(n)=\vct v(x),\qquad k<x<k+1.
\]
Then
\begin{equation}\label{eq:weighted-negativity}
 \sum_{c_n<0}|c_n|h_k(n)\ge (a-2^k)(3^{k+1}-b)\ge1.
\end{equation}
If the support is contained in $\{0,\dots,N\}$ and $H_{k,N}=\max_{0\le n\le N}h_k(n)$, then the total negative mass $\eta=\sum_{c_n<0}|c_n|$ obeys
\[
 \eta\ge \frac{1}{H_{k,N}}.
\]
\end{theorem}
\begin{proof}
Apply the functional $h_k$ to the representation.  Its value on the target is the negative integer in \eqref{eq:margin}; every term with $c_n\ge0$ contributes nonnegatively.  Moving the negative terms to the other side gives \eqref{eq:weighted-negativity}.  The final assertion follows by bounding every $h_k(n)$ by $H_{k,N}$.
\end{proof}

The weighted estimate is more informative than an unweighted one.  Tiny negative coefficients placed at enormous exponents can carry a large separator weight, so no degree-independent lower bound for ordinary negative mass should be expected without additional growth control.

\section{A cubic diagnostic interpolant}\label{sec:cubic}

The unique cubic polynomial satisfying
\[
 P(1)=1,\qquad P(2)=a,\qquad P(3)=b,\qquad P(6)=ab
\]
has the Lagrange form
\[
 P(z)=-\frac{(z-2)(z-3)(z-6)}{10}
 +a\frac{(z-1)(z-3)(z-6)}4
 -b\frac{(z-1)(z-2)(z-6)}6
 +ab\frac{(z-1)(z-2)(z-3)}{60}.
\]
Writing $P(z)=c_0+c_1z+c_2z^2+c_3z^3$ gives
\begin{align*}
 c_0&=\frac{36-45a+20b-ab}{10},\\
 c_1&=\frac{-216+405a-200b+11ab}{60},\\
 c_2&=\frac{11-25a+15b-ab}{10},\\
 c_3&=\frac{ab+15a-10b-6}{60}.
\end{align*}
Every coefficient lies in $60^{-1}\Z$.  By \cref{cor:positive-rigidity}, if all four coefficients were nonnegative, then $P$ would have to be a monomial; hence $x\in\{0,1,2,3\}$.  In every other case at least one coefficient is at most $-1/60$.

This cubic is not intended as a proof device by itself.  Its value is diagnostic: it gives the lowest-dimensional signed representation, an explicit initial point for linear programming, and a concrete instance of the negativity tax in \cref{thm:negativity-budget}.

\section{A two-moment Jensen rigidity theorem}\label{sec:jensen}

The fourth moment $M(6)$ is not actually needed if one is willing to use the transcendental constant $\alpha$ explicitly.

\begin{theorem}[Discrete Jensen rigidity]\label{thm:jensen}
Let $(p_n)$ be a probability distribution on $\N$.  Then
\begin{equation}\label{eq:jensen}
 \sum_np_n3^n\ge\left(\sum_np_n2^n\right)^\alpha.
\end{equation}
Equality holds if and only if the distribution is a point mass.
\end{theorem}
\begin{proof}
Put $U=2^N$.  Since $3^N=U^\alpha$ and $u\mapsto u^\alpha$ is strictly convex for $\alpha>1$, Jensen's inequality gives \eqref{eq:jensen}, with equality exactly when $U$, and hence $N$, is constant.
\end{proof}

Thus a positive representation matching only the values at $1,2,3$ would already prove the conjecture.  The advantage of the four-node covariance identity is arithmetic: it contains only integer powers and yields integral or rational defects.  Jensen's inequality supplies the geometry; the covariance identity supplies the exact algebraic certificate.

For a fixed integer $A\in(2^k,2^{k+1})$, the minimum possible value of $\E 3^N$ among integer-valued $N$ satisfying $\E2^N=A$ is obtained by a mixture of $k$ and $k+1$.  It is the chord value
\begin{equation}\label{eq:chord}
 L_k(A)=3^k+\frac{2\,3^k}{2^k}(A-2^k).
\end{equation}
The strict gap $L_k(A)-A^\alpha$ is another expression of discreteness.  Multiplying by $2^k$ produces the integer determinant developed next.

\section{The integer chord cocycle}\label{sec:cocycle}

Write $x=k+\theta$ with $k=\lfloor x\rfloor$ and $0<\theta<1$.  Define
\[
 \phi(\theta)=2^{\theta+1}-3^\theta-1.
\]
The graph $u\mapsto u^\alpha$ is strictly convex, so it lies below its chord from $(1,1)$ to $(2,3)$.  Equivalently,
\begin{equation}\label{eq:phi-positive}
 \phi(\theta)>0\qquad(0<\theta<1),
\end{equation}
with $\phi(0)=\phi(1)=0$.

\begin{theorem}[Integer chord determinant]\label{thm:integer-cocycle}
If $a=2^x$ and $b=3^x$ are integers and $x\notin\Z$, then
\begin{align}
 E_1
 &:=2\,3^k(a-2^k)-2^k(b-3^k)\label{eq:E1-int}\\
 &=2\,3^ka-2^kb-6^k\label{eq:E1-expanded}\\
 &=6^k\phi(\theta)\in\Z_{\ge1}.\label{eq:E1-phi}
\end{align}
More generally, for every $n\ge1$, put
\[
 m_n=\lfloor nx\rfloor,\qquad \theta_n=\{nx\}.
\]
Then $\theta_n\ne0$ and
\begin{equation}\label{eq:En}
 E_n:=2\,3^{m_n}a^n-2^{m_n}b^n-6^{m_n}
 =6^{m_n}\phi(\theta_n)\in\Z_{\ge1}.
\end{equation}
\end{theorem}
\begin{proof}
The first identity is a rearrangement.  Substituting $a=2^k2^\theta$ and $b=3^k3^\theta$ gives \eqref{eq:E1-phi}, and positivity follows from \eqref{eq:phi-positive}.  For the iterated statement use $a^n=2^{m_n}2^{\theta_n}$ and $b^n=3^{m_n}3^{\theta_n}$.  If $\theta_n=0$, then $nx\in\Z$, so $x\in\Q$ and \cref{prop:basic-cases} forces $x\in\Z$, contrary to hypothesis.
\end{proof}

The cocycle has exact local signatures in the primitive normalization.

\begin{corollary}[Primitive congruence signatures]\label{cor:cocycle-valuations}
Assume \eqref{eq:primitive-condition}.  For every $n$ with $m_n\ge2$:
\begin{enumerate}[label=(\alph*)]
\item if $a$ is odd, then $v_2(E_n)=1$;
\item if $3\nmid b$, then $v_3(E_n)=0$.
\end{enumerate}
\end{corollary}
\begin{proof}
In \eqref{eq:En}, if $a$ is odd, the first term has exact $2$-adic valuation $1$, whereas the other two have valuation at least $m_n\ge2$.  If $3\nmid b$, the middle term is a $3$-adic unit while the first and third are divisible by $3^{m_n}$.
\end{proof}

There are also immediate endpoint exclusions.  Since $a^n$ is an integer strictly larger than $2^{m_n}$ and $b^n$ is an integer strictly smaller than $3^{m_n+1}$,
\begin{align}
 \theta_n&\ge \log_2(1+2^{-m_n}),\label{eq:endpoint-low}\\
 1-\theta_n&\ge-\log_3(1-3^{-m_n-1}).\label{eq:endpoint-high}
\end{align}
These are exponentially small lower bounds.  By contrast, linear-form estimates applied to $n\log a-m\log2$ give polynomial lower bounds in $n$ for a fixed $a$ \cite{BakerWustholz1993,Matveev2000}.  Hence \eqref{eq:endpoint-low}--\eqref{eq:endpoint-high} are not, by themselves, competitive.  The potential value of $E_n$ lies in combining its Archimedean formula with the exact valuations in \cref{cor:cocycle-valuations}, not merely in using $E_n\ge1$.

\section{Geometric Newton interpolation on powers of two}\label{sec:qnewton}

Set
\[
 f(z)=z^\alpha,\qquad \alpha=\log_2 3,
\]
so that $f(2^j)=3^j$.  Let $P_N$ be the Newton interpolation polynomial through the nodes $1,2,4,\dots,2^N$:
\begin{equation}\label{eq:PN}
 P_N(z)=\sum_{n=0}^N c_n\prod_{j=0}^{n-1}(z-2^j),
\end{equation}
where an empty product is $1$.

\begin{theorem}[Closed geometric-Newton coefficients]\label{thm:qnewton-coeff}
The coefficients in \eqref{eq:PN} are
\begin{equation}\label{eq:cn-product}
 c_0=1,\qquad
 c_n=\frac{\displaystyle\prod_{j=0}^{n-1}(3-2^j)}
 {\displaystyle\prod_{j=0}^{n-1}(2^n-2^j)}\quad(n\ge1).
\end{equation}
For $n\ge2$,
\begin{equation}\label{eq:cn-sign}
 \operatorname{sgn}(c_n)=(-1)^{n-2}.
\end{equation}
\end{theorem}
\begin{proof}
The coefficient $c_n$ is the divided difference $[1,2,\dots,2^n]f$.  More generally, if $q>1$, $r>0$, and $g(q^j)=r^j$, induction in the Newton recursion gives
\[
 [1,q,\dots,q^n]g=rac{\prod_{j=0}^{n-1}(r-q^j)}{\prod_{j=0}^{n-1}(q^n-q^j)}.
\]
Taking $(q,r)=(2,3)$ gives \eqref{eq:cn-product}.  The denominator is positive.  In the numerator the factors for $j=0,1$ are positive and all factors for $j\ge2$ are negative, proving \eqref{eq:cn-sign}.
\end{proof}

\begin{table}[htbp]\centering
\caption{Exact geometric--Newton coefficients $c_n$ for the data $2^j\mapsto3^j$.}
\label{tab:qnewton-coeff}
\begin{tabular}{r r r r}
\toprule
$n$ & sign & $|\operatorname{num}(c_n)|$ & $\operatorname{den}(c_n)$\\
\midrule
0 & $+$ & 1 & 1\\
1 & $+$ & 2 & 1\\
2 & $+$ & 1 & 3\\
3 & $-$ & 1 & 84\\
4 & $+$ & 1 & 2016\\
5 & $-$ & 13 & 999936\\
6 & $+$ & 377 & 2015870976\\
7 & $-$ & 22997 & 16384999292928\\
8 & $+$ & 574925 & 106961275384233984\\
9 & $-$ & 145456025 & 13992246200663952850944\\
\bottomrule
\end{tabular}
\end{table}


The sign alternation in the coefficient sequence is canceled by the sign of the geometric Newton basis beyond the interval containing an integer argument.

\begin{theorem}[One-signed geometric tail]\label{thm:one-signed-tail}
Let $A$ be an integer with $2^k<A<2^{k+1}$.  Then every Newton summand
\[
 c_n\prod_{j=0}^{n-1}(A-2^j),\qquad n\ge k+1,
\]
has the same sign, namely $(-1)^{k-1}$.
\end{theorem}
\begin{proof}
For $n=k+1$ all basis factors are positive and \eqref{eq:cn-sign} gives sign $(-1)^{k-1}$.  For $n\ge k+2$, exactly $n-k-1$ basis factors are negative.  Hence the total sign is
\[
 (-1)^{n-2}(-1)^{n-k-1}=(-1)^{2n-k-3}=(-1)^{k-1}.
\]
\end{proof}

\subsection{A canonical entire interpolant}

Let $(u;q)_n=\prod_{j=0}^{n-1}(1-uq^j)$ denote the $q$-Pochhammer symbol.  A direct factorization of \eqref{eq:cn-product} gives
\begin{equation}\label{eq:qpoch-term}
 c_n\prod_{j=0}^{n-1}(z-2^j)
 =2^{-n}\frac{(3;1/2)_n(z;1/2)_n}{(1/2;1/2)_n}.
\end{equation}
Consequently the series
\begin{equation}\label{eq:H}
 H(z)=\sum_{n=0}^{\infty}2^{-n}
 \frac{(3;1/2)_n(z;1/2)_n}{(1/2;1/2)_n}
\end{equation}
converges uniformly on compact subsets of $\C$ and defines an entire function.  Indeed, the finite products in the quotient are uniformly bounded on compact $z$-sets, while $2^{-n}$ supplies geometric decay.  If $z=2^m$, then $(z;1/2)_n=0$ for $n>m$, so the series terminates and Newton interpolation gives
\begin{equation}\label{eq:H-interp}
 H(2^m)=3^m\qquad(m\in\N).
\end{equation}

This construction must not be confused with a proof that $H(z)=z^\alpha$.  Values on the unbounded discrete set $\{2^m\}$ do not determine an entire function without a sufficiently restrictive growth theorem.  In fact, $H$ should be regarded as a canonical $q$-hypergeometric competitor to the nonentire principal-branch function $z^\alpha$.

The one-signed tail gives monotone rational approximations to $H(A)$ for every integer $A$ between consecutive powers of two.  A potentially decisive theorem would compare $H(A)$ and $A^\alpha$ at the scale of one integer cell.

\begin{problem}[Unit-cell comparison problem]\label{prob:unit-cell}
Establish an explicit sign and a bound of the form
\[
 0<|H(A)-A^\alpha|<\dist(H(A),\Z)
\]
with the appropriate one-sided interpretation for every integer $A$ that is not a power of two.
\end{problem}

Such a result would immediately exclude $A^\alpha\in\Z$.  The formulation is intentionally concrete: both $H(A)$ and its truncation error are given by rapidly convergent rational series.

\subsection{The denominator wall}

The unreduced denominator in \eqref{eq:cn-product} factors as
\begin{equation}\label{eq:den-factor}
 \prod_{j=0}^{n-1}(2^n-2^j)
 =2^{n(n-1)/2}\prod_{\ell=1}^n(2^\ell-1).
\end{equation}
Its binary logarithm is $\Theta(n^2)$.  The analytic tail in \eqref{eq:H} decays only geometrically, as $O_A(2^{-n})$.  Therefore the naive plan ``truncate, clear all denominators, and obtain a nonzero integer smaller than one'' fails dramatically: denominator clearing loses much more than geometric convergence gains.  Any arithmetic use of \eqref{eq:H} must exploit cancellation in reduced denominators, congruences among the factors $2^\ell-1$, or a Pad\'e/Hermite construction with substantially better arithmetic efficiency.  Basic hypergeometric methods provide a natural language for this refinement \cite{GasperRahman2004}.

\section{Signed interpolation as an optimization problem}\label{sec:optimization}

For a degree bound $N$, consider the linear system
\begin{equation}\label{eq:linear-system}
 \sum_{n=0}^Nc_n=1,\quad
 \sum_{n=0}^Nc_n2^n=a,\quad
 \sum_{n=0}^Nc_n3^n=b,\quad
 \sum_{n=0}^Nc_n6^n=ab.
\end{equation}
A natural optimization is
\[
 \eta_N(a,b)=\min\left\{\sum_{n=0}^N(-c_n)_+: (c_n)\text{ satisfies \eqref{eq:linear-system}}\right\}.
\]
The positive-rigidity theorem says $\eta_N(a,b)>0$ for every noninteger candidate.  The dual separator gives the explicit lower bound
\[
 \eta_N(a,b)\ge\frac{(a-2^k)(3^{k+1}-b)}{H_{k,N}}.
\]
This linear program is attractive for two reasons.  First, all data are integers and every vertex solution is rational with denominators controlled by minors of the matrix in \eqref{eq:linear-system}.  Second, dual solutions are exponential polynomials in $n$, making them amenable to exact sign certification by Chebyshev-system techniques.

The main caveat is equally clear: $H_{k,N}$ grows like $6^N$, so the elementary dual bound decays rapidly.  A successful program needs a family of dual certificates whose coefficient norm grows much more slowly while retaining a unit integral margin at the target.  This is a precise auxiliary-function problem rather than a generic appeal to linear programming.

\section{Renormalized iterates and why derivative extraction is too weak}\label{sec:renorm}

The powers $a^n=2^{nx}$ and $b^n=3^{nx}$ provide infinitely many rational points on the normalized curve.  With $m_n=\lfloor nx\rfloor$ and $\theta_n=\{nx\}$, put
\[
 U_n=\frac{a^n}{2^{m_n}}=2^{\theta_n},\qquad
 V_n=\frac{b^n}{3^{m_n}}=3^{\theta_n}.
\]
Whenever $\theta_n$ is small, the rational number
\begin{equation}\label{eq:Rn}
 R_n=\frac{V_n-1}{U_n-1}
 =\frac{2^{m_n}(b^n-3^{m_n})}{3^{m_n}(a^n-2^{m_n})}
\end{equation}
approximates $\alpha$:
\begin{equation}\label{eq:Rn-expansion}
 R_n=\alpha+\frac{\log3\,(\log3-\log2)}{2\log2}\,\theta_n+O(\theta_n^2).
\end{equation}
This follows by dividing the Taylor expansions of $e^{\theta\log3}-1$ and $e^{\theta\log2}-1$.

Dirichlet approximation gives infinitely many $n$ with $\theta_n\ll1/n$.  But the denominator in \eqref{eq:Rn} is exponential in $m_n\asymp n$, while the error in \eqref{eq:Rn-expansion} is merely $O(1/n)$.  Expressed in terms of the rational denominator $Q_n$, the quality is only logarithmic, roughly $O(1/\log Q_n)$, vastly weaker than what an irrationality-measure contradiction requires.  Higher-order Richardson combinations replace $O(\theta_n)$ by $O(\theta_n^r)$ for fixed $r$, but this remains a power of $1/\log Q_n$.

This audit matters because derivative extraction from near-identity elements of the multiplicative group $\langle2,3\rangle$ initially looks very promising.  The calculation shows exactly why it does not compete with linear forms in logarithms: normalization creates exponentially large denominators.

\section{Other routes tested and their present obstruction}\label{sec:no-go}

\subsection{Continued fractions of $\alpha$}
For a convergent $p/q$ to $\alpha$, both
\[
 3^q-2^p\quad\text{and}\quad b^q-a^p
\]
are nonzero integers when $x$ is nonintegral.  Their logarithmic discrepancies are proportional:
\[
 q\log b-p\log a=x(q\log3-p\log2).
\]
However, the integer differences themselves are of size roughly $2^p|q\log3-p\log2|$ and $a^p|x(q\log3-p\log2)|$.  Ordinary continued-fraction accuracy is polynomial in $q$, while the prefactors are exponential.  No ``integer between zero and one'' is produced.

\subsection{Finite differences on $3$-smooth nodes}
The values $s^x$ are integers for every $s=2^m3^n$.  One can form divided differences and exponential-polynomial combinations on such nodes.  Under a common multiplicative scaling, however, a fixed integer linear combination scales as the positive power $T^x$, not as the derivative-sized quantity $T^{x-r}$.  Dividing by the Vandermonde scale destroys integrality.  A successful finite-difference argument would need unusually close additive configurations of smooth numbers together with divisibility that survives normalization; the standard multiplicative grid does not provide this.

\subsection{Positive approximation without exact interpolation}
Bernstein polynomials and other positivity-preserving schemes approximate $z^x$ very well on compact intervals.  Yet the target vector lies strictly outside the discrete positive moment cone, with the explicit distance lower bound in \cref{cor:distance}.  Approximation alone cannot cross the separator; the arithmetic mesh must force exactness before positivity rigidity becomes available.

\subsection{The canonical $q$-interpolant}
The entire function $H$ in \eqref{eq:H} is arithmetically structured and rapidly computable, but equality $H(A)=A^\alpha$ is false in general and cannot be inferred from agreement on powers of two.  The useful target is not equality; it is the one-integer-cell comparison in \cref{prob:unit-cell}.

\section{Computational reconnaissance}\label{sec:computation}

A numerical scan was performed for every integer
\[
 2^k<A<2^{k+1},\qquad 1\le k\le23,
\]
seeking the smallest distance from $A^\alpha$ to an integer in each dyadic block.  A binary64 vectorized pass selected the candidate in each block; that candidate was then reevaluated with $120$ decimal digits.  This is exploratory evidence, not a proof: a rigorous exhaustive certificate would require directed interval arithmetic for every value or an analytic lower bound.

\begin{longtable}{r r r r r}
\toprule
$k$ & minimizing $A$ & nearest $B$ & $|A^\alpha-B|$ & $\{\log_2 A\}$\\
\midrule
\endfirsthead
\toprule
$k$ & minimizing $A$ & nearest $B$ & $|A^\alpha-B|$ & $\{\log_2 A\}$\\
\midrule
\endhead
1 & 3 & 6 & \num{2.9547751e-1} & \num{5.849625e-1}\\
2 & 6 & 17 & \num{1.1356748e-1} & \num{5.849625e-1}\\
3 & 15 & 73 & \num{1.2410154e-1} & \num{9.068906e-1}\\
4 & 24 & 154 & \num{2.2107357e-2} & \num{5.849625e-1}\\
5 & 63 & 711 & \num{2.8971681e-2} & \num{9.7727992e-1}\\
6 & 102 & 1526 & \num{4.9809728e-3} & \num{6.7242534e-1}\\
7 & 163 & 3208 & \num{2.8016483e-4} & \num{3.4872815e-1}\\
8 & 331 & 9859 & \num{2.6760048e-4} & \num{3.7068741e-1}\\
9 & 589 & 24577 & \num{2.5850153e-4} & \num{2.0212382e-1}\\
10 & 1537 & 112398 & \num{3.8544369e-4} & \num{5.8590145e-1}\\
11 & 2762 & 284584 & \num{6.0204244e-4} & \num{4.314976e-1}\\
12 & 5513 & 851059 & \num{7.0582573e-5} & \num{4.2862189e-1}\\
13 & 9329 & 1959024 & \num{1.6530819e-5} & \num{1.8750673e-1}\\
14 & 26983 & 10546585 & \num{3.5403609e-5} & \num{7.1976314e-1}\\
15 & 33566 & 14906687 & \num{2.2296214e-5} & \num{3.4713003e-2}\\
16 & 125362 & 120337724 & \num{8.289363e-6} & \num{9.3574058e-1}\\
17 & 189427 & 231498127 & \num{4.724862e-7} & \num{5.3128245e-1}\\
18 & 378854 & 694494381 & \num{1.4174586e-6} & \num{5.3128245e-1}\\
19 & 566702 & 1314775264 & \num{2.1222605e-6} & \num{1.1223077e-1}\\
20 & 1302553 & 4917240094 & \num{1.3656507e-5} & \num{3.1291065e-1}\\
21 & 2774041 & 16296423255 & \num{3.9182644e-6} & \num{4.0355768e-1}\\
22 & 4332389 & 33034225344 & \num{6.7526983e-5} & \num{4.6731356e-2}\\
23 & 8431833 & 94913206833 & \num{2.6841518e-4} & \num{7.4148627e-3}\\
\bottomrule
\end{longtable}


The table serves two purposes.  It records the relevant scale of near misses for testing auxiliary functions, and it prevents a proposed inequality from being tuned only to generic points.  Any unit-cell theorem based on the canonical interpolant $H$ should be stress-tested first at these record candidates.

\section{Four focused next targets}\label{sec:targets}

\subsection{Target A: a unit-cell theorem for the $q$-interpolant}
Prove a one-sided comparison between $H(A)$ and $A^\alpha$ that is stronger than their distance to the nearest integer.  The one-signed Newton tail and the explicit $q$-hypergeometric form make this a sharply posed analytic-arithmetic problem.

\subsection{Target B: an arithmetic lower bound for weighted negativity}
Strengthen \eqref{eq:weighted-negativity} by showing that a rational signed interpolant of controlled height cannot place all of its negative mass at exponents where $h_k(n)$ is huge.  A bound coupling coefficient denominators, support height, and separator weight could turn the unit margin into a contradiction.

\subsection{Target C: an adelic obstruction for the chord cocycle}
Combine the Archimedean identity $E_n=6^{m_n}\phi(\theta_n)$ with the primitive valuations $v_2(E_n)=1$ or $v_3(E_n)=0$.  The desired theorem would rule out an integer sequence with these fixed local signatures arising from the irrational-rotation cocycle for every $n$.

\subsection{Target D: a machine-checked positivity core}
Formalize \cref{thm:covariance-identity,prop:stability,thm:separator,thm:negativity-budget}.  These results are elementary enough to make a compact trusted kernel.  A formal development would also support exact linear-program certificates produced by computation.

\section{Lean formalization blueprint}\label{sec:lean}

The following is a proof architecture, not a claim that the code has been compiled.  For finite support, define
\[
 F_c(t)=\sum_{i\in S}c_i t^i.
\]
The key algebraic lemma is the identity
\[
 F_c(1)F_c(6)-F_c(2)F_c(3)
 =\sum_{\substack{i,j\in S\\i<j}}c_ic_j(2^j-2^i)(3^j-3^i).
\]
A Lean development can proceed by expanding products with \texttt{Finset.sum\_mul}, splitting the double sum into diagonal and ordered pairs, and invoking positivity of natural powers.

\begin{lstlisting}[language={}] 
-- Schematic Lean 4 / mathlib structure; identifiers may need adjustment.
def eval (S : Finset ℕ) (c : ℕ → ℝ) (t : ℝ) : ℝ :=
  ∑ n in S, c n * t ^ n

theorem covariance_identity
    (S : Finset ℕ) (c : ℕ → ℝ) :
    eval S c 1 * eval S c 6 - eval S c 2 * eval S c 3 =
      ∑ n in S, ∑ m in S.filter (n < ·),
        c n * c m * ((2 : ℝ)^m - 2^n) * ((3 : ℝ)^m - 3^n) := by
  -- expand, pair (n,m) with (m,n), cancel diagonal
  sorry

theorem support_subsingleton_of_eq_zero
    (S : Finset ℕ) (c : ℕ → ℝ)
    (hc : ∀ n ∈ S, 0 ≤ c n)
    (hzero : eval S c 1 * eval S c 6 = eval S c 2 * eval S c 3) :
    (S.filter (fun n => c n ≠ 0)).card ≤ 1 := by
  rw [← sub_eq_zero] at hzero
  rw [covariance_identity] at hzero
  -- every summand is nonnegative; a positive pair is impossible
  sorry
\end{lstlisting}

The infinite-series theorem can then be obtained either by monotone convergence or by truncating and passing to the limit.  The separator theorem is even shorter: prove the sign of two monotone factors according to whether $n\le k$ or $k+1\le n$.

\section{Conclusions}\label{sec:conclusion}

The work in this report does not prove the Alaoglu--Erd\H{o}s conjecture.  It does produce a new, compact rigidity mechanism and several exact consequences:
\begin{enumerate}
\item The four-node defect is a sum of strictly positive pair interactions.  Zero defect forces a single integer exponent.
\item Near-zero defect quantitatively concentrates any positive representing measure, and rational weights give a discrete defect gap.
\item A four-term exponential polynomial separates every noninteger exponent from the integer moment cone with an integral margin.
\item Every signed interpolant must pay at least one unit of separator-weighted negative mass.
\item Geometric Newton coefficients for $2^j\mapsto3^j$ have a closed product formula and produce a one-signed tail at integer arguments.
\item Every primitive hypothetical counterexample yields a positive integer cocycle with exact $2$-adic or $3$-adic signatures.
\end{enumerate}

The positivity theorem is strong enough that the remaining issue can no longer be described merely as ``find a better interpolant.''  A successful proof must extract positivity, or a substitute carrying the same equality rigidity, from arithmetic data that naturally produce signed objects.  Of the routes developed here, the most concrete immediate target is the unit-cell comparison for the canonical $q$-interpolant; the most structural is an arithmetic strengthening of the weighted-negativity budget.

\appendix
\section{Proof of the general geometric divided-difference formula}\label{app:qproof}

For completeness, let $q>1$, let $r$ be a scalar, and let $G_n(z)$ be the Newton polynomial of degree at most $n$ satisfying $G_n(q^m)=r^m$ for $0\le m\le n$.  Write
\[
 G_n(z)=\sum_{j=0}^n d_j\prod_{\ell=0}^{j-1}(z-q^\ell).
\]
We prove
\[
 d_n=\frac{\prod_{\ell=0}^{n-1}(r-q^\ell)}{\prod_{\ell=0}^{n-1}(q^n-q^\ell)}.
\]
For $n=0$ this is $d_0=1$.  Evaluating the Newton expansion at $z=q^n$ gives
\[
 r^n-G_{n-1}(q^n)=d_n\prod_{\ell=0}^{n-1}(q^n-q^\ell).
\]
The finite $q$-binomial identity
\[
 G_{n-1}(q^n)=r^n-\prod_{\ell=0}^{n-1}(r-q^\ell)
\]
follows by induction on $n$: both sides are polynomials of degree $n$ in $r$, agree at $r=q^j$ for $0\le j<n$, and have the same leading coefficient.  Substitution gives the claimed formula.

\section{Reproducible numerical scan}\label{app:code}

The following compact script reproduces the dyadic-block scan in \cref{sec:computation}.  It is intentionally labeled experimental rather than rigorous.

\begin{lstlisting}[language=Python]
import numpy as np
import mpmath as mp

mp.mp.dps = 120
alpha_mp = mp.log(3) / mp.log(2)
alpha = float(alpha_mp)

for k in range(1, 24):
    lo, hi = 2**k + 1, 2**(k+1)
    best_A, best_d = None, 1.0
    for s in range(lo, hi, 1_000_000):
        e = min(s + 1_000_000, hi)
        A = np.arange(s, e, dtype=np.float64)
        Y = np.exp(alpha * np.log(A))
        D = np.abs(Y - np.rint(Y))
        j = int(np.argmin(D))
        if float(D[j]) < best_d:
            best_d = float(D[j])
            best_A = s + j
    y = mp.power(mp.mpf(best_A), alpha_mp)
    nearest = mp.nint(y)
    print(k, best_A, nearest, abs(y-nearest))
\end{lstlisting}

\begin{thebibliography}{99}
\bibitem{AlaogluErdos1944}
L.~Alaoglu and P.~Erd\H{o}s,
\emph{On highly composite and similar numbers},
Transactions of the American Mathematical Society \textbf{56} (1944), 448--469.

\bibitem{Ax1971}
J.~Ax,
\emph{On Schanuel's conjectures},
Annals of Mathematics \textbf{93} (1971), 252--268.

\bibitem{Baker1975}
A.~Baker,
\emph{Transcendental Number Theory},
Cambridge University Press, 1975.

\bibitem{BakerWustholz1993}
A.~Baker and G.~W\"ustholz,
\emph{Logarithmic forms and group varieties},
Journal f\"ur die reine und angewandte Mathematik \textbf{442} (1993), 19--62.

\bibitem{GasperRahman2004}
G.~Gasper and M.~Rahman,
\emph{Basic Hypergeometric Series}, second edition,
Cambridge University Press, 2004.

\bibitem{Gelfond1934}
A.~O.~Gelfond,
\emph{Sur le septi\`eme probl\`eme de Hilbert},
Bulletin de l'Acad\'emie des Sciences de l'URSS (1934), 623--634.

\bibitem{HardyLittlewoodPolya1952}
G.~H.~Hardy, J.~E.~Littlewood, and G.~P\'olya,
\emph{Inequalities}, second edition,
Cambridge University Press, 1952.

\bibitem{Karlin1968}
S.~Karlin,
\emph{Total Positivity, Vol. I},
Stanford University Press, 1968.

\bibitem{Lang1966}
S.~Lang,
\emph{Introduction to Transcendental Numbers},
Addison--Wesley, 1966.

\bibitem{LaurentMignotteNesterenko1995}
M.~Laurent, M.~Mignotte, and Y.~Nesterenko,
\emph{Formes lin\'eaires en deux logarithmes et d\'eterminants d'interpolation},
Journal of Number Theory \textbf{55} (1995), 285--321.

\bibitem{Matveev2000}
E.~M.~Matveev,
\emph{An explicit lower bound for a homogeneous rational linear form in logarithms of algebraic numbers. II},
Izvestiya: Mathematics \textbf{64} (2000), 1217--1269.

\bibitem{Schneider1934}
T.~Schneider,
\emph{Transzendenzuntersuchungen periodischer Funktionen. I, II},
Journal f\"ur die reine und angewandte Mathematik \textbf{172} (1934), 65--69, 70--74.

\bibitem{Waldschmidt2000}
M.~Waldschmidt,
\emph{Diophantine Approximation on Linear Algebraic Groups},
Grundlehren der mathematischen Wissenschaften 326, Springer, 2000.
\end{thebibliography}

\end{document}
